\documentclass[12pt,a4paper]{ctexbook}
\usepackage{geometry}
\geometry{top=2.5cm,bottom=2.5cm,left=2.5cm,right=2.5cm}
\usepackage{amsmath,amssymb,mathtools}
\usepackage{enumitem,array,booktabs,pifont}
\usepackage{hyperref,fancyhdr,bookmark}

\pagestyle{fancy}
\fancyhf{}
\fancyhead[C]{\leftmark}
\fancyfoot[C]{\thepage}
\title{\textbf{数学 · 备考指南}}
\date{}

\begin{document}
\maketitle
\tableofcontents

\chapter{答题时间整体规划与答题技巧}

\section{试卷结构速览}
\begin{center}
\begin{tabular}{c|c|c|c}
题型 & 题量 & 分值 & 建议用时 \\ \hline
单选题 & 8 题 & 40 分 & 25 分钟 \\
多选题 & 4 题 & 20 分 & 15 分钟 \\
填空题 & 4 题 & 20 分 & 15 分钟 \\
解答题 17 题 & 1 题 & 10 分 & 10 分钟 \\
解答题 18--20 题 & 3 题 & 36 分 & 30 分钟 \\
解答题 21--22 题 & 2 题 & 24 分 & 25 分钟 \\ \hline
总计 & 22 题 & 150 分 & 120 分钟
\end{tabular}
\end{center}

\section{答题顺序策略}
\textbf{原则}：先易后难，先熟后生，遇难则跳。\\
\textbf{推荐顺序}：单选1-6 → 填空13-14 → 解答17-18 → 多选9-10 → 解答19-20 → 单选7-8 → 多选11-12 → 填空15-16 → 解答21 → 解答22 → 检查

\section{选择题技巧}
\textbf{直接法}：公式直接套用；\textbf{特值法}：取$0,\pm1,\pm2$快速验证；\textbf{排除法}：奇偶性/正负号/定义域排除；\textbf{数形结合}：零点画图、方程根分布；\textbf{估算法}：不精确计算，只比大小

\section{多选题策略}
核心：\textbf{宁少勿滥}。只选100\%确定的选项，不确定的宁愿不选（2分 vs 丢5分的取舍）

\section{填空题注意}
无过程分，注意：定义域、多解、$\pm$符号、化简结果、区间开闭

\section{大题踩分要点}
\textbf{三角}：公式+代入+结果；\textbf{数列}：通项推导+求和；\textbf{立几}：建系+坐标+法向量；\textbf{概统}：分布列+期望公式；\textbf{解几}：联立+韦达+条件转化；\textbf{导数}：求导+讨论+极值

\chapter{高考考频分析}

\section{近五年考点分布}
\begin{center}
\begin{tabular}{c|c|c|c|c|c}
模块 & 2021 & 2022 & 2023 & 2024 & 2025 \\ \hline
集合 & 5 & 5 & 5 & 5 & 5 \\
复数 & 5 & 5 & 5 & 5 & 5 \\
函数与导数 & 27 & 27 & 32 & 32 & 27 \\
三角与解三角形 & 10 & 10 & 17 & 15 & 10 \\
数列 & 12 & 17 & 10 & 10 & 12 \\
向量 & 5 & 5 & 0 & 5 & 5 \\
立体几何 & 17 & 12 & 17 & 12 & 12 \\
解析几何 & 17 & 22 & 17 & 17 & 22 \\
概率统计 & 17 & 17 & 17 & 12 & 17 \\
排列组合 & 5 & 5 & 5 & 5 & 5
\end{tabular}
\end{center}

\section{必考模块分值}
\begin{itemize}
  \item 函数与导数：25-32分（选填1-2道+压轴解答）
  \item 解析几何：17-22分（选填1道+解答21题）
  \item 立体几何：12-17分（选填1道+解答19题）
  \item 概率统计：12-17分（选填1道+解答20题）
  \item 三角、数列：10-17分（交替出现）
  \item 集合、复数、排列组合：各5分（选填）
\end{itemize}

\section{命题趋势}
\begin{enumerate}
  \item 反套路化——三角和数列交替出现在17/18题
  \item 情境题增多——概率统计结合实际情境
  \item 导数难度增加——隐零点、同构、极值点偏移
  \item 解析几何淡化计算——强化代数转化能力
  \item 新定义题型出现——考察现场学习能力
\end{enumerate}

\chapter{真题代练}

\section{集合}
\textbf{例1}（2024Ⅰ卷）$A=\{x\mid -5<x^3<5\}$，$B=\{-3,-1,0,2,3\}$，$A\cap B=$\\
解析：$\sqrt[3]{5}\approx1.71$，$A\approx(-1.71,1.71)$，$A\cap B=\{-1,0\}$ \quad 选A

\section{复数}
\textbf{例2} $z=\frac{1-i}{2+2i}$，求$z-\bar{z}$\\
解析：$z=-\frac12i$，$\bar{z}=\frac12i$，$z-\bar{z}=-i$ \quad 选A

\section{函数与导数}
\textbf{例3} $f(x)=\begin{cases}e^x+a,&x\le0\\ \ln(x+1)+x,&x>0\end{cases}$在$\mathbb{R}$上递增，求$a$范围\\
解析：$a\le0$（衔接处），$a\ge-1$（值域连续性），$\therefore a\in[-1,0]$

\textbf{例4} $a=0.1e^{0.1}$，$b=\frac19$，$c=-\ln0.9$，比大小\\
解析：$a\approx0.1105$，$b\approx0.1111$，$c\approx0.1054$，$c<a<b$

\section{三角}
\textbf{例5} $\triangle ABC$中$A+B=3C$，$2\sin(A-C)=\sin B$，$AB=5$，求$\sin A$和$AB$边上高\\
解析：$C=45^\circ$，$\tan A=3$，$\sin A=\frac{3\sqrt{10}}{10}$，高$=6$

\section{数列}
\textbf{例6} $a_1=1$，$a_{n+1}=a_n+1$（$n$奇）或$a_n+3$（$n$偶），求$b_n=a_{2n}$及$S_{20}$\\
解析：$b_n=4n-2$，$S_{20}=300$

\section{概率统计}
\textbf{例7} 甲命中0.6，乙命中0.8，命则继续，未中则换人。第1次甲乙各0.5，求第2次为乙概率\\
解析：$P=0.4\times0.5+0.8\times0.5=0.6$

\section{解析几何}
\textbf{例8} 椭圆$\frac{x^2}{a^2}+\frac{y^2}{b^2}=1$右焦点$F$，$M(1,\frac32)$在$E$上且$MF\perp x$轴，过$F$直线交$E$于$A,B$，$N$为$AB$中点，求$ON$斜率最大值\\
解析：$E:\frac{x^2}{4}+\frac{y^2}{3}=1$，$k_{\max}=\frac{\sqrt3}{4}$

\chapter{分类训练}

\section{三角恒等变换}
\begin{enumerate}
  \item[★] $\sin\alpha=\frac35$，$\alpha\in(\frac{\pi}{2},\pi)$，求$\cos2\alpha$和$\tan(\alpha+\frac{\pi}{4})$
  \item[★] $\triangle ABC$中$a=7,b=8,\cos B=-\frac17$，求$\angle A$
  \item[★★] $\triangle ABC$中$2b\cos C=2a-c$，求$B$；若$b=2\sqrt3$，求面积最大值
\end{enumerate}

\section{数列}
\begin{enumerate}
  \item[★] 等差$\{a_n\}$中$a_1=1,a_3=5$，求$S_{10}$
  \item[★] 等比$\{a_n\}$中$a_1=2,a_4=16$，求$a_n$和$S_5$
  \item[★★] $a_1=1$，$a_{n+1}=2a_n+1$，求$a_n$
\end{enumerate}

\section{立体几何}
\begin{enumerate}
  \item[★] 正方体中求$A_1B$与$B_1C$所成角
  \item[★★] 四棱锥$P-ABCD$中底面矩形，$PA\perp$底面，$PA=AD=2,AB=3$，$E$为$PD$中点，求二面角$E-AC-B$余弦值
\end{enumerate}

\section{概率统计}
\begin{enumerate}
  \item[★] 3红2白取2球，求恰1红概率和红球数期望
  \item[★★] 甲胜率0.6乙胜率0.4，三局两胜，求甲获胜概率
\end{enumerate}

\section{解析几何}
\begin{enumerate}
  \item[★] 椭圆$\frac{x^2}{4}+\frac{y^2}{3}=1$，$F$为右焦点，过$F$直线交椭圆于$A,B$，求$\triangle AOB$面积最大值
  \item[★★] 双曲线$\frac{x^2}{3}-y^2=1$右焦点$F$，过$F$直线交双曲线于$A,B$，若$\overrightarrow{OA}\cdot\overrightarrow{OB}=1$，求$AB$方程
\end{enumerate}

\section{函数与导数}
\begin{enumerate}
  \item[★] $f(x)=x^3-3x^2+1$的单调区间和极值
  \item[★★] $f(x)=\ln x - a(x-1)$，讨论单调性
  \item[★★★] $f(x)=e^x-ax-1$，若$f(x)\ge0$恒成立，求$a$
\end{enumerate}

\appendix
\chapter{分类训练答案}
\textbf{三角1.} $\cos2\alpha=2(-\frac45)^2-1=\frac{7}{25}$，$\tan(\alpha+\frac{\pi}{4})=\frac{-\frac34+1}{1-(-\frac34)}=\frac{\frac14}{\frac74}=\frac17$\\
\textbf{三角2.} $\sin B=\frac{4\sqrt3}{7}$，$\frac{7}{\sin A}=\frac{8}{4\sqrt3/7}\Rightarrow\sin A=\frac{\sqrt3}{2}$，$\because b>a\therefore A=60^\circ$\\
\textbf{三角3.} $B=60^\circ$，$S_{\max}=3$\\
\textbf{数列1.} $d=2$，$S_{10}=100$\\
\textbf{数列2.} $q=2$，$a_n=2^n$，$S_5=62$\\
\textbf{数列3.} $a_n=2^n-1$\\
\textbf{立几1.} $60^\circ$\\
\textbf{立几2.} $\cos\theta=\frac{3}{\sqrt{13}}$\\
\textbf{概统1.} $P=\frac35$，$E(X)=1.2$\\
\textbf{概统2.} $P=0.6^2+2\cdot0.6^2\cdot0.4=0.648$\\
\textbf{解几1.} $S_{\max}=\sqrt3$\\
\textbf{解几2.} $x=\pm\sqrt3 y+2$\\
\textbf{导数1.} 增$(-\infty,0),(2,+\infty)$减$(0,2)$，极$f(0)=1,f(2)=-3$\\
\textbf{导数2.} $a\le0$时递增；$a>0$时$(0,\frac1a)$增$(\frac1a,+\infty)$减\\
\textbf{导数3.} $a=1$

\end{document}
